% 第 4 周周记（2026.08.24 - 08.28）
\documentclass[UTF8,fontset=fandol]{ctexart}
% =====================================================================
% internship/preamble.tex
% 实习周记统一模板的共用导言区。各周文件在 \documentclass 之后用
%   % =====================================================================
% internship/preamble.tex
% 实习周记统一模板的共用导言区。各周文件在 \documentclass 之后用
%   % =====================================================================
% internship/preamble.tex
% 实习周记统一模板的共用导言区。各周文件在 \documentclass 之后用
%   \input{../preamble.tex}
% 引入本文件，正文前用 \weekhead{周次}{时间}{主题} 输出抬头。
% 编译命令：xelatex weekN.tex
% =====================================================================

\usepackage[a4paper,top=2.6cm,bottom=2.6cm,left=2.8cm,right=2.8cm]{geometry}
\usepackage{booktabs}   % 表格横线，week3 的结果对比表会用到
\usepackage{array}

% 周记抬头：\weekhead{周次}{时间区间}{本周主题}
\newcommand{\weekhead}[3]{%
\begin{center}
{\zihao{3}\bfseries 实习周记（#1）}\\[10pt]
{\zihao{-4}姓名：【姓名】\qquad 学号：【学号】}\\[4pt]
{\zihao{-4}时间：#2}\\[4pt]
{\zihao{-4}主题：#3}
\end{center}
\vspace{2pt}
\hrule height 0.8pt
\vspace{10pt}
\zihao{-4}%
}

% 引入本文件，正文前用 \weekhead{周次}{时间}{主题} 输出抬头。
% 编译命令：xelatex weekN.tex
% =====================================================================

\usepackage[a4paper,top=2.6cm,bottom=2.6cm,left=2.8cm,right=2.8cm]{geometry}
\usepackage{booktabs}   % 表格横线，week3 的结果对比表会用到
\usepackage{array}

% 周记抬头：\weekhead{周次}{时间区间}{本周主题}
\newcommand{\weekhead}[3]{%
\begin{center}
{\zihao{3}\bfseries 实习周记（#1）}\\[10pt]
{\zihao{-4}姓名：【姓名】\qquad 学号：【学号】}\\[4pt]
{\zihao{-4}时间：#2}\\[4pt]
{\zihao{-4}主题：#3}
\end{center}
\vspace{2pt}
\hrule height 0.8pt
\vspace{10pt}
\zihao{-4}%
}

% 引入本文件，正文前用 \weekhead{周次}{时间}{主题} 输出抬头。
% 编译命令：xelatex weekN.tex
% =====================================================================

\usepackage[a4paper,top=2.6cm,bottom=2.6cm,left=2.8cm,right=2.8cm]{geometry}
\usepackage{booktabs}   % 表格横线，week3 的结果对比表会用到
\usepackage{array}

% 周记抬头：\weekhead{周次}{时间区间}{本周主题}
\newcommand{\weekhead}[3]{%
\begin{center}
{\zihao{3}\bfseries 实习周记（#1）}\\[10pt]
{\zihao{-4}姓名：【姓名】\qquad 学号：【学号】}\\[4pt]
{\zihao{-4}时间：#2}\\[4pt]
{\zihao{-4}主题：#3}
\end{center}
\vspace{2pt}
\hrule height 0.8pt
\vspace{10pt}
\zihao{-4}%
}


\begin{document}

\weekhead{第 4 周}{2026 年 8 月 24 日至 8 月 28 日}{论文写作与修改、发明专利申请文件撰写}

实验收口之后，本周的重心转到写作和专利。论文按与导师商定的结构组织：引言、相关工作、预备知识与问题形式化、系统设计、评估。写英文初稿最花时间的是把方案表述成严格的定义：pattern key 是什么，历史并集如何定义，recall 和 precision 的公式怎么写，优化问题如何形式化。动笔前我先建了一张术语表，把 online/offline path、pattern key、training window 等词的用法固定下来，后面十几轮修改全靠它保持一致，否则每改一版都会冒出新旧词混用的问题。

初稿完成后的修改量超出我的预期，三轮意见累计一百多条，多数是表述和数据问题。最痛的一处是正文里写的"超过 80\% 的 pattern key，其历史并集小于 10 个槽"，导师在批注里说这句话和配图对不上，怀疑是模型产生的幻觉数据。我重新统计后发现，实际是 83.5\% 的键并集小于 50 个槽。写作时凭印象写数据被当场抓包，这个教训我记到现在，之后的每个数字都从 CSV 或代码输出里核对。其它有代表性的修改包括：块数口径要全文统一，采集到的 1,001 个块过滤后剩 999 个，两个口径的差异要在文中解释；正文和算法里同一个符号出现了两种写法，需要统一；算法排版用的宏包相互冲突，行号一度全部显示为 0，调试了半天；所有图按统一的蓝灰色系重画，曲线加标记，去掉图内文字，只留 LaTeX 的图注。

论文定稿后，同一套方案写成发明专利。我此前没有接触过专利申请，按导师提供的写作指南，参照多份真实授权专利逆向出句式模板，按说明书、附图、权利要求书、摘要的顺序逐项完成。摘要要求不超过 300 字，改到约 280 字。附图共 4 张，做了中英文两套版本。写权利要求时最需要适应的是回到"必要技术特征"的思维，独权里只放绕不开的步骤，实现细节和效果数据都留到说明书里，技术效果还要给出量化数据支撑创造性。和论文相比，专利面向的读者和使用的语言都完全不同。

周末把四周的工作整理成周记和实习报告。回头看，这条线的脉络是先被一组真实的延迟数字说服，再用数据画像找到访问规律，接着用一个负面的可行性结论把方案逼到正确的方向，最后在写作中把每一步讲清楚。过程中犯过的错，口径不清、凭印象写数据、实验规模不够，最后都变成了写作时的检查清单。四周时间很紧，如果重来一次，我会把论文大纲和术语表更早定下来，写作和实验并行推进，而不是等所有结果都跑完再动笔。

\end{document}
